\documentclass[a4paper, oneside]{article}
\usepackage{graphicx}
\usepackage{amsthm}
\usepackage{amsmath}
\usepackage{amssymb}
\usepackage[a4paper,
            bindingoffset=0.2in,
            left=2cm,
            right=2cm,
            top=2cm,
            bottom=2cm,
            footskip=.25in]{geometry}
\usepackage[italian]{babel}
\usepackage{pgfplots}
\usepackage{tabularx}
\usepackage{tikz}
\usepackage{wrapfig}
\usepackage{color}
\usepackage[d]{esvect}
\definecolor{page}{rgb}{0.129,0.157,0.212}
\pagecolor{page}
\color{white}
\graphicspath{ {./images/} }
\usetikzlibrary{shapes.geometric}
\usetikzlibrary{datavisualization}
\usetikzlibrary{datavisualization.formats.functions}
\usetikzlibrary{patterns}
\pgfplotsset{width=10cm,compat=1.18}

\title{Appunti di Metodi}
\author{Tommaso Miliani}
\date{11-03-26}

\begin{document}
\newtheoremstyle{theoremEnv}
                {}          % Space above
                {}          % Space below
                {\slshape}  % Body font
                {}          % Indent amount
                {\bfseries} % Head font
                {.}         % Punctuation after head
                {\newline}  % Space after theorem head
                {}          % Theorem head spec
\theoremstyle{theoremEnv}

\newtheorem{definition}{Definizione}[section]
\newtheorem{theorem}{Teorema}[section]
\newtheorem{lemma}{Proposizione}[section]
\newtheorem{observation}{Osservazione}[section]
\newtheorem{corollary}{Corollario}[theorem]
\newtheorem{example}{Esempio}[section]
\newtheorem{remark}{Enunciato}[section]

\maketitle

\section{Funzioni derivabili in senso complesso (olomorfe)}
Si può riscrivere l'identità di Cauchy-Riemann come
\begin{gather*}
    \partial x f = \partial x u + i \partial x v \\
    \partial y f = \partial y u + i \partial y v
\end{gather*}
Per cui
\begin{gather*}
    0 = \partial x u - \partial y v + (\partial x v + \partial y u)
\end{gather*}
Queste due parti (sia la reale che l'immaginaria) devono essere zero, 
ottenendo una forma equivalente dell'identità:
\begin{gather*}
    \left\{\begin{array}{l}
        \partial x u = \partial y v \\
        \partial x v = - \partial y u
    \end{array}\right.
\end{gather*}
Che prende il nome di identità di Cauchy-Riemann in forma reale.

\begin{theorem}
    Se $f$ è una funzione complessa e derivabile in $z = x + iy$ senso complesso, 
    allora valgono le identità di Cauchy-Riemann. \\
    Se valgono le identità di Cauchy-Riemann in un certo punto e
    $u$ e $v$ sono continue in $z = x + iy$, allora $f = u + iv$ è
    derivabile in senso complesso in $z$. 
\end{theorem}

\begin{definition}[Operatore di Cauchy-Riemann]
    Preso un numero complesso $z = x + iy$, si può considerare $\overline{z} = z^{\star} = x - iy$.
    Queste due variabili devono essere interpretate come indipendenti  
    \begin{gather*}
        \left\{\begin{array}{l}
            z = x + iy \\
            \overline{z} = x + iy 
        \end{array}\right. \ \Longrightarrow \ \left\{\begin{array}{l}
            x = \dfrac{1}{2}(z + \overline{z} ) \\
            y = \dfrac{1}{2i}(z - \overline{z} )
        \end{array}\right.
    \end{gather*}
    Formalmente si interpretano come indipendenti: se si avesse una funzione $F(x, y)$, essa
    sarebbe scrivibile come $\mathcal{F}(z, \overline{z} )$, scrivendo $x$ e $y$
    come visto sopra. Si introducono due operatori differenziali
    \begin{align}
        \partial z &= \frac{\partial}{\partial z} = \frac{1}{2}(\partial x - i\partial y) \\
        \partial \overline{z} &= \frac{\partial }{\partial \overline{z} } = \frac{1}{2} (\partial x + i \partial y)  
    \end{align}
    Che prendono il nome di \textbf{operatori di Cauchy-Riemann}. Qualsiasi funzione olomorfa soddisfa
    la seguente condizione
    \begin{align}
        \partial \overline{z} f(z, \overline{z} ) = \frac{1}{2}(\partial f + i \partial y f) = 0 
    \end{align}
    Se questo vale, $f$ non è più funzione di $\overline{z}$ ma solo di $z$. Analogamente se
    $\partial z f(z, \overline{z}) = 0$, $f$ diventa funzione solo di $\overline{z}$ e prende il 
    nome di \textbf{funzione anti-olomorfa} per cui valgono tutte le proprietà delle funzioni 
    olomorfe se si trasforma $\overline{z}$ in $z$. È una scrittura che 
    mima un cambio di variabile ma non è assolutamente un cambio di variabile. 
\end{definition}

\section{Derivabilità su di un aperto}
\begin{definition}
    Presa una funzione $f$ definita su di un aperto con $\mathbb{A}$ un aperto su $\mathbb{C}$. $f$ è dunque 
    olomorfa se è derivabile in senso complesso e $f'$ è continua. Si definisce
    $H(\mathbb{A})$ come l'insieme delle funzioni olomorfe sull'aperto $\mathbb{A}$. 
\end{definition} \noindent
SI è già visto come la derivata in senso complesso dei polinomi sia 
la stessa della derivata reale. In realtà tutte le regole rimangono uguali se la funzione è
olomorfa:
\begin{itemize}
    \item $(f + g)' = f' + g'$
    \item $(fg)' = f'g + fg'$
    \item $\left(\dfrac{f}{g}\right)' = (f'g - fg') / g^{2}  $
    \item $ \left(f(g(z))\right)' = (f \circ g)' = (f' \circ g) g'$
    \item $(f^{-1})^{-1} = \dfrac{1}{f' \circ f^{-1}}$
\end{itemize}
Le funzioni $u$ e $v$ che si ottengono a partire dalle parti reali e immaginarie,
hanno delle proprietà interessanti se si utilizza l'identità di Cauchu-Riemann:
\begin{gather*}
    f = u + iv \\
    \partial^{2}x u = \partial x \partial y v = \partial y \partial x v = - \partial ^{2}y u
\end{gather*}
Lo scambio dell'ordine della derivata si può fare perché si utilizza il teorema di Schwartz 
(si scambiano se esistono e se sono continue). Allora 
\begin{gather*}
    (\partial^{2} x + \partial^{2}y)u = \Delta u = 0 
\end{gather*}
Dunque $u$ è una \textbf{funzione armonica}. Si fa lo stesso per $v$:
\begin{gather*}
    (\partial^{2}x + \partial^{2} y) v = \Delta v = 0
\end{gather*}
Le due funzioni sono \textbf{coniugate} l'una rispetto all'altra perché
sono legate dall'identità di Cauchy-Riemann, lo stesso vale per le 
funzioni anti-olomorfe. Inoltre
\begin{gather*}
    \Delta = 4\partial z\partial \overline{z} = 4\partial \overline{z} \partial z  
\end{gather*}
Dove $\Delta$ è il \textbf{operatore Laplaciano}.

\section{Integrazione di funzioni complesse}
L'integrale di una funzione complessa sfrutta un integrale di linea 
poiché si può interpretare $\mathbb{C}$ come spazio bidimensionale: 
si può dunque integrare sul cammino di entrambe le variabili. Si definisce
\begin{gather*}
    u(x) = f(x) \ dx 
\end{gather*}
Se si volesse integrare questa funzione da $\gamma : [a, b] \to \mathbb{R}$:
\begin{gather*}
    \int_{\gamma} = \int_{\overline{a} }^{\overline{b} }  f(x) \ dx = \int_{a}^{b} f(\gamma(t)) \gamma'(t) \ dt 
\end{gather*}
Se siamo in $n$ dimensioni si utilizza solo l'ultimo integrale per cui 
$\gamma : [a, b] \to \mathbb{R}^{n}$. Su $\mathbb{C}$ si vuole fare lo stesso, 
ma interpretando $\gamma$ come se fosse un cammino che va dal reale al complesso.

\begin{definition}
    Sia $\gamma : [a, b] \to \mathbb{A} \subset \mathbb{C}$ e $f : \mathbb{A} \to \mathbb{C}$, con $\mathbb{A}$
    un aperto. Allora si definisce l'integrale
    \begin{align}
        \int_{\gamma} f(z) \ dz = \int_{a}^{b} f(\gamma(t))\gamma'(t) \ dt
    \end{align}
    Come l'integrale di cammino su $\mathbb{C}$ (ed è orientato).
    In questo caso $f$ è una uno-forma su $\mathbb{C}$ o una uno-forma su $\mathbb{R}^{2}$.
    DUnque 
    \begin{gather*}
        \int_{\gamma}f = -\int_{-\gamma } f
    \end{gather*}
\end{definition}

\begin{example}
    Preso $\gamma: [0, 1]$ che manda $t \to z_1 + (z_2 - z_1)t $, che è 
    un segmento che va da $z_1$ a $z_2$:
    \begin{gather*}
        \int_{\gamma} 1 \ dz = \int_{0}^{1} (z_2 - z_1) \ dt = z_2 - z_1 \\
        \int_{\gamma} z \ dx = \int_{0}^{1} (z_1 +(z_2 - z_1) t) (z_2 - z_1) \ dt = z_1(z_2 - z_1)t + \frac{1}{2}(z_2 - z_1)^{2}    
    \end{gather*}
\end{example}

\begin{example}
    Preso un cerchio ri raggio $r$ attorno all'origine:

    Il cammino $\gamma : [0, 2\pi]$ e manda $\theta \to z_0 + re^{i\theta}$. Si vuole fare
    \begin{gather*}
        \oint_{\gamma} (z - z_0)^{n} \ dz 
    \end{gather*}
    lungo $\gamma$. Allora 
    \begin{gather*}
        \int_{0}^{2pi} ((z_0 + re^{i\theta}) - z_0)^{n} \frac{1}{d\theta}(z_0 + re^{i\theta})\ d\theta = \int_{0}^{2\pi} r^{n} e^{i\theta} re^{i\theta}i\ d\theta = ir^{n + 1}\int_{0}^{2\pi} e^{i(n + 1)1\theta} \ d\theta    
    \end{gather*}
    Se $n = -1$ fa $2\pi i$, altrimenti fa zero. Infatti il risultato è $2\pi i \delta_{n - 1}M$, dove $M$ indica il numero di 
    avvolgimenti che compie la curva intorno a $z_0$. 
\end{example}


\begin{theorem}[Teorema di Cauchy]
    Presa $f$ una funzione olomorfa su $\mathbb{A}$ un aperto semplicemente
    connesso (ogni curva può essere deformato con continuità fino ad arrivare ad un punto), 
    per ogni $\gamma \subset \mathbb{A}$, 
    \begin{align}
        \oint_{\gamma} f(z) \ dz = 0 
    \end{align}
    Dunque si ha integrale nullo su di un cammino. 
\end{theorem}
\begin{proof}
    \begin{gather*}
        \oint _{\gamma} f(z ) \ dz = \oint _\gamma (u + iv)(dx + i dy)
    \end{gather*}
    In questo modo posso ottenere un normale integrale di linea su $\mathbb{C}$:
    \begin{gather*}
        \oint_{\gamma}^{} (u dx - v dy)  + \oint_\gamma  (v\ dx + u\ dy)
    \end{gather*}
    Dunque sono due integrali da $\mathbb{R}^{2}$ ad $\mathbb{R}$. Con il teorema di Gauss-Green
    si ottiene che 
    \begin{gather*}
        \oint_{\gamma = \partial \Gamma} u dx - u dy = \int\int_{\Gamma} (-\partial x v - \partial y u) \ dx \ dy  = 0
    \end{gather*}
    Dove $\Gamma \subset \mathbb{A}$ ed è la porzione del dominio racchiusa da $\gamma$, dunque
    $\gamma$ è la frontiera di $\Gamma$. 
    \begin{gather*}
        \oint _{\gamma = \partial \Gamma} (v \ dx + u dy ) = \int \int_\Gamma (\partial x u - \partial y v) dx dy = 0 
    \end{gather*}
\end{proof}

\begin{example}[Esempio di applicazione su di un dominio non nullomotopo]
    L'integrale sul circuito che circonda $z_0$:
    \begin{gather*}
        \oint_\gamma \frac{1}{z - z_0} = 2\pi i
    \end{gather*}
    Il dominio che circonda $z_0$ non è nullomotopo in quanto 
    non è semplicemente connesso dunque l'intergrale non fa zero. Se invece
    si facesse lo stesso integrale su di un circuito che non circondi $z_0$, allora può fare zero. 
\end{example}

\begin{theorem}
    Sia $f$ olomorfa su di un insieme $\mathbb{A}$. Se $\beta, \gamma$ sono circuiti 
    \textbf{omotopi} su $\mathbb{A}$, allora 
    \begin{align}
        \oint_{\beta} f = \oint_\gamma  f
    \end{align}
    Connettendono i due circuiti, anche se ci fosse un buco nel mezzo, si ottiene un circuito nullomotopo, infatti
    \begin{align}
        \oint_\beta f +\oint_{-\gamma}f = 0
    \end{align}
\end{theorem}

\begin{theorem}
    Posti $\beta$ e $\gamma$ dei cammini omotopi in un aperto $\mathbb{A}$ con 
    stessa origine e stessa estremità finale, e presa una funzione
    $f$ olomorfa allora 
    \begin{align}
        \int_{\beta} f = \int_{\gamma} f  
    \end{align}
\end{theorem}
\begin{proof}
    \begin{gather*}
        0 = \int_{\beta}^{} f - \int_{\gamma}^{} f = \int_{\beta}^{} f + \int_{-\gamma}^{}   f
    \end{gather*}
    Dato che $\beta$ e $\gamma$ sono omotopi, allora la somma dei due circuiti è nullomotopa, 
    per il teorema di Cauchy allora 
    \begin{gather*}
        \int_{\beta - \gamma}^{} f = 0 
    \end{gather*}
    Ossia la tesi. 
\end{proof}

\begin{definition}[Primitiva di una funzione]
    Sia $F$ una funzione olomorfa su di un aperto $\mathbb{A}$, si definisce 
    $F$ la \textbf{primitiva} di $f$ su $\mathbb{C}$ se
    \begin{align}
        F'(z) = f(z) \qquad \forall z \in \mathbb{A}
    \end{align}
\end{definition}

\begin{theorem}[Teorema della primitiva]
    Sia $f : \mathbb{A} \to \mathbb{C}$ continua sull'aperto $\mathbb{A}$. Allora
    le seguenti affermazioni sono equivalenti:
    \begin{enumerate}
        \item $f$ ammette primitiva
        \item Se $\beta, \gamma$ sono cammini omotopi $ \subset \mathbb{A}$ 
        con le stesse estremità, allora 
        \begin{align}
            \int_{\beta}f = \int_{\gamma}^{} f  
        \end{align}
        \item $\forall \gamma \subset  \mathbb{A}$ si ha che 
        \begin{align}
            \oint_\gamma f = 0 
        \end{align}
    \end{enumerate}
    Si dimostra solo che $1 \ \Longrightarrow \ 2$.
\end{theorem}
\begin{proof}
    Si assume che esista $F' = f$, con $F \in H(\mathbb{A})$, allora
    per definizione si ha che $\gamma : [a, b] \to \mathbb{C}$:
    \begin{gather*}
        \int_{\gamma}^{} f dz = \int_{a}^{b}  F'(\gamma(t))\gamma'(t) \ dt = \int_{a}^{b} (F \circ \gamma)' (t) \ dt  
    \end{gather*}
    Dato che è la derivata composta, 
    \begin{gather*}
        = F(\gamma(b)) - F(\gamma(a))
    \end{gather*}
    Ossia la differenza tra il punto iniziale e quello finale. 
\end{proof}

\begin{definition}[Lunghezza del cammino]
    Posta $\gamma$ una curva su di un aperto generico, la lunghezza del cammino 
    è data dalla seguente:
    \begin{align}
        L_\gamma \equiv \int_{a}^{b} \left| \gamma'(t) \right| \ dt  
    \end{align}
\end{definition}
Si introduce ora il seguente teorema:
\begin{theorem}[Disuguaglianza di Darboux]
    \begin{align}
        \left| \int_{\gamma} f \ dz  \right| \leq \sup_{\gamma} \left| f \right| L_\gamma   
    \end{align}
\end{theorem}
\begin{proof}
    L'integrale di una funzione reale a valori complessi è 
    \begin{gather*}
        \left| \int_{a}^{b} g(t) \ dt  \right| \leq \int_{a}^{b} \left| g(t) \right| \ dt    
    \end{gather*}
    Questo vale anche se la funzione integranda ha valori complessi (ossia se $g(t) \in \mathbb{C}$).
    \begin{gather*}
        \left| \int_{\gamma}^{} f \right| = \left| \int_{a}^{b} f(\gamma(t)) \gamma'(t)   \ dt \right| \leq  \int_{a}^{b} \left| f(\gamma(t)) \right| \left| \gamma'(t) \right|  dt \leq \sup_{\gamma}\left| f \right| \int_{a}^{b} \left| \gamma'(t) \right| \ dt = \sup_{\gamma}\left| f \right| L_\gamma    
    \end{gather*}
\end{proof}


\section{Funzioni analitiche su $\mathbb{C}$}
\begin{theorem}[formulaintegreale di Cauchy]
    Sia $f \in H(\mathbb{A})$ sull'aperto $\mathbb{A}$. SIa $z \in \mathbb{A}$ e 
    sia $\gamma$ un circuito nullomotopo su $\mathbb{A}$ che circonda $z$ una volta
    sola in verso positivo. Allora
    \begin{align}
        f(z) = \frac{1}{2\pi i} \int_{\gamma} \frac{f(w)}{w - z} dw 
    \end{align}
    Questo teorema mi permette, conoscendo la funzione sul bordo del circuit, di conoscere la funzione in tutti i punti 
    interni del circuito.
\end{theorem}
\begin{proof}
    Dimostrazione sulle dispense (da fare!) 
\end{proof}






\end{document}